\documentclass[conference]{IEEEtran}
\usepackage{amsmath}
\usepackage{graphicx}

\title{Sparse Retrieval Still Wins at Equal Cost}
\author{A. Researcher}

\begin{document}
\maketitle

\begin{abstract}
Retrieval quality on this benchmark has plateaued since 2021. We show that
BM25 with query expansion matches a dense encoder at equal cost.
\end{abstract}

\section{Introduction}
\label{sec:intro}

Dense encoders are reported to beat sparse baselines \cite{smith2021,nakamura2023},
but the comparisons in Fig.~\ref{fig:budget} hold the index size constant
rather than the budget. % TODO: soften this before submission

% This whole line is a note to my co-author and is not part of the paper.
The gap closes once cost is held fixed, as \citet{chen2024} also observed.

\begin{figure}[t]
  \centering
  \includegraphics[width=\columnwidth]{budget.pdf}
  \caption{Recall against budget. This caption is not a claim sentence.}
  \label{fig:budget}
\end{figure}

\section{Method}

The corpus is indexed with BM25 using $k_1 = 0.9$ and $b = 0.4$. Queries are
expanded with the top three terms from a first-pass retrieval, which follows
the setup in Sec.~3 of the original paper.

\begin{equation}
  \label{eq:score}
  s(q, d) = \sum_{t \in q} \mathrm{idf}(t) \cdot \frac{f(t, d)}{f(t, d) + k_1}
\end{equation}

The scoring function above is unchanged from the reference implementation.


\section{Results}

Recall at ten improves by four points at one third of the cost, i.e. the gain
is not an artifact of one query set. The inline term $k_1 = 0.9$ was fixed
throughout, and Eq.~\ref{eq:score} is unchanged.

\begin{table}[t]
  \caption{Systems compared.}
  \begin{tabular}{lrr}
    \toprule
    System & Recall@10 & Cost \\
    \midrule
    BM25+QE & 0.71 & 1.0 \\
    Dense & 0.69 & 3.2 \\
    \bottomrule
  \end{tabular}
\end{table}

Three conditions were tested.

\begin{itemize}
  \item Sparse retrieval alone, with no expansion applied to the query.
  \item Sparse retrieval with the expansion described in \emph{Method}.
  \item A dense encoder trained on the same corpus.
\end{itemize}

\bibliographystyle{IEEEtran}
\bibliography{refs}

\end{document}
